\section{2009级工科数学分析下 A卷}

\subsection{资料来源}
\begin{itemize}
  \item 试题：2009级计算机工科数学分析下 期末A.doc，DOC 正文已抽取；公式对象已按 Word 公式图片逐项恢复。
  \item 公式图片审计：\texttt{tmp/past-exam-equations/2009-A/manifest.csv}。
\end{itemize}

\subsection{试题}
诚信应考，考试作弊将带来严重后果！

华南理工大学 2009--2010 第二学期期末考试

《工科数学分析》（下）试卷（A）

注意事项：
\begin{enumerate}
  \item 考前请将密封线内填写清楚；
  \item 所有答案请直接答在试卷上；
  \item 考试形式：闭卷；
  \item 本试卷共五大题，满分 100 分，考试时间 120 分钟。
\end{enumerate}

\subsubsection*{一、填空题（共 5 小题，每小题 4 分，共 20 分）}
\begin{enumerate}
  \item $(y''')^3+(y')^4+x^2+x+1=0$ 是 \underline{\hspace{2cm}} 阶微分方程。
  \item 已知 $(a\times b)\cdot c=2$，则 $[(a+b)\times(b+c)]\cdot(c+a)=$ \underline{\hspace{3cm}}。
  \item 函数 $z=\arctan\dfrac{y}{x}+\ln\sqrt{x^2+y^2}$ 在点 $(1,0)$ 处的梯度 $\operatorname{grad}z\big|_{(1,0)}=$ \underline{\hspace{3cm}}。
  \item 交换积分顺序，则 $\displaystyle\int_1^e dx\int_0^{\ln x}f(x,y)\,dy=$ \underline{\hspace{3cm}}。
  \item 设 $f(x)$ 是以 $2\pi$ 为周期的函数，它在 $[-\pi,\pi]$ 上的表达式是
  \[
  f(x)=
  \begin{cases}
  x+1,&-\pi\le x<0,\\
  x^2,&0\le x<\pi.
  \end{cases}
  \]
  若它的傅里叶级数的和函数为 $S(x)$，则 $S(0)=$ \underline{\hspace{3cm}}。
\end{enumerate}

\subsubsection*{二、选择题（共 5 小题，每小题 4 分，共 20 分）}
\begin{enumerate}
  \item 抛物线 $y^2=4x$ 上与直线 $x-y+4=0$ 相距最近的点是（\quad）。

  A. $(0,0)$；\quad B. $(1,1)$；\quad C. $(1,2)$；\quad D. $(2,2\sqrt2)$。

  \item 直线 $L:\dfrac{x-2}{2}=\dfrac{y+1}{-2}=\dfrac{z-3}{1}$ 与平面 $\pi:x+2y-2z=6$ 的关系是（\quad）。

  A. 平行；\quad B. 垂直；\quad C. 相交但不垂直；\quad D. 重合。

  \item 设 $f(x,y)$ 在点 $(x_0,y_0)$ 处的两个偏导数都存在，则（\quad）。

  A. $f(x,y)$ 在点 $(x_0,y_0)$ 必连续；\quad
  B. $f(x,y)$ 在点 $(x_0,y_0)$ 必可微；\quad
  C. $\displaystyle\lim_{x\to x_0}f(x,y_0)$ 与 $\displaystyle\lim_{y\to y_0}f(x_0,y)$ 都存在；\quad
  D. $\displaystyle\lim_{\substack{x\to x_0\\y\to y_0}}f(x,y)$ 存在。

  \item 设区域 $D:0\le y\le x,\ 0\le x\le\pi$，则二重积分 $\displaystyle\iint_D\sqrt{1-\sin^2x}\,dxdy$（\quad）。

  A. $0$；\quad B. $\dfrac{\pi}{2}$；\quad C. $2$；\quad D. $\pi$。

  \item 设 $\alpha$ 为常数，则级数 $\displaystyle\sum_{n=1}^{\infty}\left[\dfrac{\sin(n\alpha)}{n^2}-\dfrac1n\right]$（\quad）。

  A. 绝对收敛；\quad B. 发散；\quad C. 条件收敛；\quad D. 收敛性与 $\alpha$ 的取值有关。
\end{enumerate}

\subsubsection*{三、计算题（本大题共 4 小题，每小题 8 分，共 32 分）}
\begin{enumerate}
  \item 求微分方程 $xy''=y'\ln\dfrac{y'}{x}$ 的通解。
  \item 计算三重积分 $\displaystyle I=\iiint_{\Omega}(x^2+y^2+z^2)\,dV$，其中 $\Omega:x^2+y^2+z^2\le 2z$。
  \item 计算曲线积分
  \[
  I=\int_L (x+y)^2\,dx-(x^2+y^2\sin y)\,dy,
  \]
  其中 $L$ 是抛物线 $y=x^2$ 上从点 $(-1,1)$ 到点 $(1,1)$ 的那一段。
  \item 求幂级数 $\displaystyle\sum_{n=0}^{\infty}(2n+1)x^{2n}$ 的收敛域及和函数。
\end{enumerate}

\subsubsection*{四、应用题（本大题共 2 小题，每小题 7 分，共 14 分）}
\begin{enumerate}
  \item 造一个容积为 $V$ 的长方形无盖铝盒，怎样设计尺寸才能使它的表面积最小？
  \item 求密度 $\rho\equiv1$ 的均匀球壳 $x^2+y^2+z^2=a^2\ (z\ge0)$ 对于 $Oz$ 轴的转动惯量。
\end{enumerate}

\subsubsection*{五、证明题（本大题共 2 小题，每小题 7 分，共 14 分）}
\begin{enumerate}
  \item 设函数 $f(x)$ 在区间 $[a,b]$ 上连续，证明：
  \[
  \int_a^b dx\int_a^x f(y)\,dy=\int_a^b f(y)(b-y)\,dy.
  \]
  \item 设正数序列 $\{x_n\}$ 单调上升且有界，证明级数
  \[
  \sum_{n=1}^{\infty}\left(1-\frac{x_n}{x_{n+1}}\right)
  \]
  收敛。
\end{enumerate}

\subsection{答案与解析}

\subsubsection*{一、填空题}
\begin{enumerate}
  \item 最高阶导数为 $y'''$，故为 $3$ 阶微分方程。
  \item 展开得
  \[
  [(a+b)\times(b+c)]\cdot(c+a)
  =(a\times b)\cdot c+(b\times c)\cdot a=2+2=4.
  \]
  \item
  \[
  z_x=\frac{x-y}{x^2+y^2},\qquad z_y=\frac{x+y}{x^2+y^2}.
  \]
  在 $(1,0)$ 处 $\operatorname{grad}z=(1,1)$。
  \item 积分区域为 $1\le x\le e,\ 0\le y\le\ln x$，即 $0\le y\le1,\ e^y\le x\le e$。因此
  \[
  \int_1^e dx\int_0^{\ln x}f(x,y)\,dy
  =\int_0^1dy\int_{e^y}^e f(x,y)\,dx.
  \]
  \item $x=0$ 是跳跃点，傅里叶级数在该点收敛到左右极限平均值：
  \[
  S(0)=\frac{f(0-)+f(0+)}2=\frac{1+0}{2}=\frac12.
  \]
\end{enumerate}

\subsubsection*{二、选择题}
\begin{enumerate}
  \item 令 $y=t,\ x=t^2/4$，到直线的距离只需最小化
  \[
  \left(\frac{t^2}{4}-t+4\right)^2.
  \]
  在选项中最小点为 $(1,2)$，故选 C。
  \item 直线方向向量 $v=(2,-2,1)$，平面法向量 $n=(1,2,-2)$。因 $v\cdot n=-4\ne0$，直线与平面相交；且 $v$ 不平行于 $n$，故不垂直于平面。选 C。
  \item 偏导数存在只保证沿坐标方向的一元极限存在，不保证连续、可微或二重极限存在。选 C。
  \item
  \[
  \iint_D\sqrt{1-\sin^2x}\,dxdy
  =\int_0^\pi x|\cos x|\,dx=\pi.
  \]
  故选 D。
  \item $\sum\sin(n\alpha)/n^2$ 绝对收敛，而 $-\sum1/n$ 发散，所以原级数发散。选 B。
\end{enumerate}

\subsubsection*{三、计算题}
\begin{enumerate}
  \item 设 $p=y'$，则
  \[
  xp'=p\ln\frac{p}{x}.
  \]
  令 $q=\ln\dfrac{p}{x}$，有 $p=xe^q$，代入得
  \[
  1+xq'=q,\qquad q=1+C_1x.
  \]
  因而 $y'=ex e^{C_1x}$。故
  \[
  y=C_2+e\int x e^{C_1x}\,dx.
  \]
  即当 $C_1\ne0$ 时
  \[
  y=C_2+\frac{e\,e^{C_1x}(C_1x-1)}{C_1^2};
  \]
  当 $C_1=0$ 时
  \[
  y=C_2+\frac e2x^2.
  \]
  \item 区域为球心 $(0,0,1)$、半径 $1$ 的球。令 $z=w+1$，则
  \[
  x^2+y^2+z^2=x^2+y^2+w^2+1+2w.
  \]
  对称性使 $2w$ 的积分为 $0$，所以
  \[
  I=\iiint_{r\le1}r^2\,dV+\iiint_{r\le1}1\,dV
  =\frac{4\pi}{5}+\frac{4\pi}{3}
  =\frac{32\pi}{15}.
  \]
  \item 取参数 $x=t,\ y=t^2,\ -1\le t\le1$，则 $dx=dt,\ dy=2t\,dt$。于是
  \[
  I=\int_{-1}^{1}\left[(t+t^2)^2-2t(t^2+t^4\sin t^2)\right]dt
  =\int_{-1}^{1}(t^2+t^4)\,dt=\frac{16}{15}.
  \]
  \item 令 $r=x^2$，则
  \[
  \sum_{n=0}^{\infty}(2n+1)x^{2n}
  =2\sum_{n=0}^{\infty}nr^n+\sum_{n=0}^{\infty}r^n
  =\frac{1+r}{(1-r)^2}.
  \]
  因此收敛域为 $(-1,1)$，和函数为
  \[
  S(x)=\frac{1+x^2}{(1-x^2)^2}.
  \]
\end{enumerate}

\subsubsection*{四、应用题}
\begin{enumerate}
  \item 设底面长宽为 $x,y$，高为 $h$，则 $xyh=V$，表面积
  \[
  S=xy+2xh+2yh.
  \]
  极小值由对称性取 $x=y=a$。于是 $h=V/a^2$，
  \[
  S=a^2+\frac{4V}{a},\qquad S'=2a-\frac{4V}{a^2}.
  \]
  得 $a^3=2V$。故底面为正方形，边长
  \[
  a=(2V)^{1/3},\qquad h=\frac a2=\left(\frac V4\right)^{1/3}.
  \]
  \item 在上半球面上取球坐标，$x^2+y^2=a^2\sin^2\varphi$，$dS=a^2\sin\varphi\,d\varphi d\theta$。于是
  \[
  I_{Oz}=\iint_\Sigma(x^2+y^2)\,dS
  =2\pi a^4\int_0^{\pi/2}\sin^3\varphi\,d\varphi
  =\frac{4\pi a^4}{3}.
  \]
\end{enumerate}

\subsubsection*{五、证明题}
\begin{enumerate}
  \item 左端积分区域为 $a\le y\le x\le b$。交换积分次序：
  \[
  \int_a^b dx\int_a^x f(y)\,dy
  =\int_a^b f(y)\,dy\int_y^b dx
  =\int_a^b f(y)(b-y)\,dy.
  \]
  \item 因 $\{x_n\}$ 为正数且单调上升，有 $x_{n+1}\ge x_1>0$。于是
  \[
  0\le1-\frac{x_n}{x_{n+1}}
  =\frac{x_{n+1}-x_n}{x_{n+1}}
  \le\frac{x_{n+1}-x_n}{x_1}.
  \]
  部分和满足
  \[
  \sum_{n=1}^N\left(1-\frac{x_n}{x_{n+1}}\right)
  \le\frac{x_{N+1}-x_1}{x_1},
  \]
  右端有界，故正项级数收敛。
\end{enumerate}
